The following template provides model language for the compactness section of
an expert report filed in redistricting litigation.
Bracketed items [\ldots] should be replaced with plan-specific values.

\bigskip
\hrule
\bigskip

\noindent\textbf{COMPACTNESS ANALYSIS} \\
\textit{[Expert Name], Expert Report, [Case Name], [Date]}

\bigskip

\noindent I have computed six standard compactness metrics for each district in
the proposed [name] congressional redistricting plan for [State], applied to
2020 Census TIGER/Line tract geometries projected to the Albers Equal Area
Conic coordinate system (EPSG:5070).
All metrics were computed using the \textsc{bisect} algorithmic redistricting
platform, version [version], with seed $= [seed]$.

\bigskip

\noindent\textbf{Mean Compactness Profile} (all [k] districts):
\begin{itemize}
  \item \textbf{Polsby-Popper (PP)}: Mean $= [PP]$ (scale $[0,1]$; higher is more compact;
    $1.0$ = circle; $0.785$ = square). Source: K.1.
  \item \textbf{Reock}: Mean $= [Reock]$ (scale $[0,1]$; higher is more compact;
    less sensitive to coastal boundary artifacts than PP). Source: K.2.
  \item \textbf{Convex Hull Ratio (CH)}: Mean $= [CH]$ (scale $(0,1]$; $1.0$
    means perfectly convex; no district has tentacle appendages below CH $= 0.85$).
    Source: K.3.
  \item \textbf{Schwartzberg (S)}: Mean $= [S]$ (scale $[1,\infty)$; lower is
    more compact; $S = 1/\sqrt{PP}$). Source: K.4.
  \item \textbf{Length-Width Ratio (LW)}: Mean $= [LW]$; maximum district LW $= [LW\_max]$.
    No district exceeds the LW $= 3$ court threshold. Source: K.5.
  \item \textbf{Population-Weighted Compactness (PWC)}: Mean $= [PWC]$ km$^2$.
    [Algorithm] achieves [X]\% lower PWC than a standard bisection baseline.
    Source: K.6.
\end{itemize}

\noindent The submitted plan achieves PP $= [PP]$, Reock $= [Reock]$,
CH $= [CH]$, placing it in the [percentile] of algorithmically-generated plans
for [State] under 2020 Census data.

\noindent All six metrics are within the acceptable range for all $[k]$ districts
(PP $\geq 0.18$, Reock $\geq 0.30$, CH $\geq 0.82$, S $\leq 2.4$, LW $\leq 2.8$,
PWC within two standard deviations of state mean).

\noindent The complete per-district compactness data are provided in Appendix~A
(machine-readable JSON) and Appendix~B (district-by-district table).

\bigskip
\hrule
\bigskip

\subsection{Notes on the Template}

\paragraph{Percentile language.}
The ``percentile'' claim requires running the \textsc{bisect} ensemble
(\texttt{--search bisection-ensemble}) to generate a distribution of plans
and report the submitted plan's percentile position.
For a single-seed plan (seed $= 42^\dagger$), the percentile is approximate.

\paragraph{Certification statement.}
For certified court filings, add:
\begin{quote}
  \emph{I certify that the compactness values reported herein were computed
  using version [X.Y.Z] of the \textsc{bisect} platform from the TIGER/Line
  2020 Census tract geometry files (SHA-256: [hash]) with no modification to
  the underlying data or software.}
\end{quote}
